%% The first command in your LaTeX source must be the \documentclass command.
\documentclass[hf]{ceurart}

\sloppy

\usepackage{listings}
\usepackage{graphicx}
\usepackage{booktabs}
\usepackage{microtype}
\usepackage{setspace}
\usepackage{xcolor}

% ---- Listings style for Cypher (UTF-8 safe) ----
\definecolor{cypherkw}{rgb}{0.0,0.0,0.6}
\lstdefinestyle{cypher}{
  basicstyle=\ttfamily\scriptsize,
  breaklines=true,
  breakatwhitespace=false,
  columns=fullflexible,
  keepspaces=true,
  showstringspaces=false,
  frame=single,
  framesep=4pt,
  extendedchars=true,
  inputencoding=utf8,
  keywordstyle=\color{cypherkw}\bfseries,
  morekeywords={MATCH,WHERE,RETURN,AND,OR,NOT,AS,WITH,toLower,CONTAINS,count,DISTINCT},
  literate=
    {á}{{\'a}}1 {à}{{\`a}}1 {â}{{\^a}}1 {ã}{{\~a}}1 {ä}{{\"a}}1
    {é}{{\'e}}1 {ê}{{\^e}}1 {è}{{\`e}}1 {í}{{\'i}}1 {ï}{{\"i}}1
    {ó}{{\'o}}1 {ô}{{\^o}}1 {õ}{{\~o}}1 {ö}{{\"o}}1
    {ú}{{\'u}}1 {ü}{{\"u}}1 {ç}{{\c c}}1 {Á}{{\'A}}1 {É}{{\'E}}1
}
\lstset{breaklines=true}

\begin{document}

\copyrightyear{2026}
\copyrightclause{Copyright for this paper by its authors.
  Use permitted under Creative Commons License Attribution 4.0
  International (CC BY 4.0).}

\conference{2nd International Workshop on Semantic Technologies for Data Management
(ST4DM 2026), July 18, 2026, Lisbon, Portugal}

\title{Using Knowledge Graphs to Restrict the Context of LLM Analysis:
A GraphRAG Approach for Personnel Management}

\author[1]{Antonio Gustavo Silveira Dantas}[
  email=antonio.dantas101835@ga.ita.br,
]
\cormark[1]
\address[1]{Núcleo Corporativo de Tecnologia da Informação,
            São José dos Campos, SP, Brazil}

\author[2]{José Maria Parente de Oliveira}[
  email=parente@ita.br
]
\address[2]{Instituto Tecnológico de Aeronáutica,
            São José dos Campos, SP, Brazil}

%-----------------------------------------------------------------------
\begin{abstract}
%-----------------------------------------------------------------------
We investigate the combination of Knowledge Graphs (KGs) and Large Language
Models (LLMs) as a means of grounding model responses in verified organizational
data and reducing hallucinations. A domain ontology covering people, skills,
positions, and projects was defined and instantiated in a Neo4j graph database.
A graph-based retrieval-augmented (GraphRAG) pipeline was built on top of this
graph: natural language questions are translated into Cypher queries, executed
against the KG, and the retrieved structured data is supplied as context for LLM
response generation. Three architectures were evaluated on a manually validated
gold-standard corpus of 95 question--query--response triples: GPT-4 standalone,
a locally-run quantized model fine-tuned for Cypher generation
(text-to-cypher-Gemma-3-27B, referred to as HF\,27B), and a hybrid combining
both. Factual consistency was judged by Claude (model \texttt{claude-opus-4-8}), an
LLM that is not one of the evaluated systems. HF\,27B achieved a perfect 100\%
Cypher execution rate against 97.9\% for GPT-4, but only 49\% factual consistency
(95\% CI 40--59) in the generated answers. GPT-4 reached 70\% (95\% CI 61--79) and
the hybrid 59\% (95\% CI 49--68). The fully local pipeline trails the fully cloud
one by a wide, significant margin (49\% vs.\ 70\%, McNemar $p<0.001$); decomposing
this gap shows that the query-generation and answer-generation stages contribute
comparably (each isolated step about +10 points, $p\approx0.06$), so a semantically
wrong query and a weak summarizer of verbose results are both liabilities and
neither can be skipped. Results support GraphRAG as a viable and practical
strategy for domains requiring factual traceability, such as strategic personnel
planning and legal research.
\end{abstract}

\begin{keywords}
  Knowledge Graphs \sep
  Large Language Models \sep
  GraphRAG \sep
  Hallucination Mitigation \sep
  Personnel Management
\end{keywords}

\maketitle

%=======================================================================
\section{Introduction}
%=======================================================================

The growing adoption of Large Language Models such as
GPT-4~\cite{openai2023gpt4} and LLaMA~\cite{touvron2023llama} across
professional settings has brought increased attention to a well-documented
failure mode: hallucination. Because these models generate text by sampling
from distributions learned during training, they tend to produce fluent but
factually unfounded statements whenever a query touches knowledge that was
underrepresented in the training corpus or requires precise symbolic
reasoning~\cite{mallen2023trustlanguagemodelsinvestigating,
huang2023reasoninglargelanguagemodels}. This is not a fringe
concern---surveys have documented hallucination rates severe enough to render
LLMs unreliable for high-stakes tasks without additional
safeguards~\cite{agrawal2024knowledgegraphsreducehallucinations}.

At the same time, Knowledge Graphs (KGs) have become a mature technology for
representing domain knowledge in a structured, navigable, and formally grounded
form~\cite{hogan2021knowledge}. An ontology provides the conceptual
backbone~\cite{GRUBER1993199}, defining which entities exist and what
relationships are valid between them; a graph database such as
Neo4j~\cite{Neo4j2024} makes those structures efficiently queryable. The core
hypothesis of this work is that pairing a KG with an LLM---using the graph as
the authoritative source of facts and the model as a natural-language
interface---can substantially reduce hallucination without sacrificing
conversational accessibility~\cite{agrawal2024knowledgegraphsreducehallucinations,
UnifyingLLMandKG}.

Retrieval-Augmented Generation (RAG)~\cite{lewis2021retrievalaugmentedgenerationknowledgeintensivenlp}
provides the architectural pattern for this integration. Standard RAG uses
dense-vector similarity search to retrieve relevant documents; graph-based RAG
replaces that with symbolic graph queries, which return exactly the facts
matching a pattern rather than approximate nearest neighbors. This is
particularly advantageous for the kind of multi-hop relational queries that
arise in organizational settings---\textit{``Which engineers with Python expertise
have worked on cybersecurity projects in the last two years and are eligible for
promotion?''}---where vector retrieval tends to fragment and
miss~\cite{edge2025localglobalgraphrag, gao2024retrievalaugmentedgenerationlargelanguage}.

The application domain we chose to study is strategic personnel management, a
field with high information complexity and clear costs for errors. Organizations,
particularly public ones with large and geographically distributed workforces,
routinely need to cross-reference skills, project histories, certifications, and
career stages to make staffing and succession decisions. This is precisely the kind
of structured, multi-entity query that KGs handle naturally and that LLMs alone
handle poorly.

This paper makes the following contributions:
\begin{enumerate}
  \item A purpose-built domain ontology and Neo4j Knowledge Graph for
        organizational personnel and knowledge management, together with an
        explicit account of its design methodology (competency questions) and a
        discussion of vocabulary reuse and the property-graph-vs-RDF design
        decision (Section~\ref{sec:ontology}).
  \item A GraphRAG pipeline evaluated under three LLM configurations: GPT-4
        standalone, a Cypher-specialized open-source model (HF\,27B), and a hybrid
        combining both, compared on a manually validated 95-question benchmark
        with confidence intervals and paired significance testing.
  \item A structured, data-driven failure analysis that separates syntactic query
        errors, semantic query errors, and answer-generation (summarization)
        errors, quantifying how often each occurs and which one binds overall
        performance.
  \item A discussion of deployment trade-offs---cost, latency, data sovereignty,
        and accuracy---that extends the results to practical organizational decisions.
\end{enumerate}

%=======================================================================
\section{Background and Related Work}
%=======================================================================

\subsection{Large Language Models and the Hallucination Problem}

LLMs are neural networks based on the Transformer
architecture~\cite{vaswani2017attention} trained on large text corpora to model
conditional token distributions~\cite{brown2020language}. At inference time,
they autoregressively generate text by sampling the most probable next token given
a context window. This mechanism makes them powerful generalists but also
fundamentally probabilistic: they do not retrieve stored facts but reconstruct
plausible continuations~\cite{mallen2023trustlanguagemodelsinvestigating}.
When training data is sparse or inconsistent for a given topic, the model may
produce confident fabrications. Injecting verified facts as context at inference
time is the main mitigation
strategy~\cite{gao2024retrievalaugmentedgenerationlargelanguage,
agrawal2024knowledgegraphsreducehallucinations}.

\subsection{Knowledge Graphs and Ontologies}
\label{sec:kg}

A graph $G(V,E)$ is a set of vertices $V$ and edges $E$ connecting pairs of
vertices~\cite{diestel2017graph}. In a \emph{labeled property graph}---the model
used by Neo4j---vertices and edges carry type labels and key-value attributes,
enabling rich, typed relational representations. A Knowledge Graph combines a
graph with a formal ontology, an explicit specification of a
conceptualization~\cite{GRUBER1993199} that defines the entity types and
admissible relationships of a domain~\cite{hogan2021knowledge}. The ontology
provides the schema against which instances are stored and queried.

Two modelling traditions coexist. The Semantic Web stack (RDF/OWL, queried with
SPARQL) encodes the ontology itself as data and supports formal reasoning and
standardized constraint languages such as SHACL. Labeled property graphs
prioritize traversal performance and developer ergonomics but, unlike RDF/OWL,
have \emph{no standardized way of encoding an ontology as part of the
graph}~\cite{hogan2021knowledge}: the schema is typically enforced by the
application layer rather than by the data model. We return to this trade-off, and
to our reasons for adopting a property graph, in Section~\ref{sec:ontology}.
It is also worth tempering a common contrast: relational databases do support
schema changes (\texttt{ALTER}) and cross-table navigation (\texttt{JOIN}); the
advantage of a property graph is not that these are impossible elsewhere but that
it \emph{facilitates} schema evolution and makes multi-hop traversal a native,
first-class operation rather than a chain of joins.

\subsection{Retrieval-Augmented Generation and GraphRAG}
\label{sec:rag}

RAG, introduced by Lewis et al.~\cite{lewis2021retrievalaugmentedgenerationknowledgeintensivenlp},
augments LLM generation with an external retrieval step: relevant documents are
fetched and prepended to the prompt so that the model answers from evidence rather
than from parametric memory alone. Standard implementations retrieve by dense
vector similarity~\cite{gao2024retrievalaugmentedgenerationlargelanguage}, which
works well for localized queries but breaks down on global ones---queries that
require synthesizing information spread across many documents or entity
chains~\cite{edge2025localglobalgraphrag}.

\emph{Graph-based RAG} replaces or augments vector retrieval with structured graph
access. In the literature the term covers a spectrum. At one end, systems such as
the GraphRAG of Edge et al.~\cite{edge2025localglobalgraphrag} build a graph from
\emph{unstructured} text using an LLM and then traverse community summaries to
answer query-focused-summarization questions. At the other end---the setting of
this paper---the graph already exists as curated structured data, and retrieval is
performed by translating the user's question into a formal graph query (here,
Cypher) whose result set is the evidence. Our pipeline therefore belongs to the
text-to-query / KG question-answering family: it is ``GraphRAG'' in the sense that
generation is augmented with facts retrieved from a graph, but it does not perform
the multi-hop neighborhood exploration or LLM-driven graph construction associated
with the unstructured variant. We make this scope explicit because the two
families have different failure modes; ours concentrates the risk in the
question-to-query translation step. The LangChain
library~\cite{LangChain2023} provides infrastructure for this pattern through its
\texttt{GraphCypherQAChain} abstraction.

\subsection{Context Selection and Related Approaches}

Even when the graph is the right source, supplying it well matters: irrelevant
graph neighborhoods inflate the prompt and degrade response quality. Several lines
of work inform our design. KGPool~\cite{nadgeri-etal-2021-kgpool} performs dynamic
context selection, conditioning graph expansion on the input sentence; the lesson
we adopt is that context should be \emph{conditioned on the question} rather than
fixed---in our pipeline the Cypher query plays exactly this role, selecting the
subgraph relevant to the question instead of dumping a neighborhood. KG-CF~\cite{KG-CF}
and the retrieve-rank-reason loop of Xu et al.~\cite{xu2024contextgraph} show that
\emph{pruning} uninformative paths improves downstream accuracy; this motivates our
choice to restrict context to the precise result set returned by the query rather
than the full neighborhood of the matched entities. At the level of integration
patterns, Pan et al.~\cite{UnifyingLLMandKG} map the landscape of LLM--KG
synergy---LLMs assisting KG construction and KGs grounding LLM inference---and our
system instantiates the second direction. The chief difference between these works
and ours is one of task: they target relation extraction or KG completion, whereas
we use the KG as a fixed, authoritative store and study the reliability of the
natural-language interface over it.

A practical concern for evaluation is benchmark quality. Zhang et
al.~\cite{zhang2025diagnosingaddressingpitfallskgrag} identified annotation
artifacts in widely used KG-RAG datasets that inflate reported accuracy and mask
real failure modes. Their findings motivated the manual validation procedure we
describe in Section~\ref{sec:dataset}.

%=======================================================================
\section{Methodology}
%=======================================================================

\subsection{Domain Ontology and Knowledge Graph Construction}
\label{sec:ontology}

\paragraph{Design methodology.}
The ontology was developed following a lightweight, competency-question-driven
process in the spirit of established ontology-engineering methodologies. We first
elicited the information needs of strategic personnel management as a set of
\emph{competency questions}---natural-language questions the system should be able
to answer, such as ``which people of a given rank hold a given skill?'', ``which
projects require a given technology?'', and ``which knowledge domains are held by
fewer than $k$ people?''. Each candidate class and relationship was admitted only
if it was required to answer at least one competency question that could not
otherwise be expressed precisely~\cite{GRUBER1993199}. We emphasize that these
competency questions guided the \emph{schema}; they are conceptually distinct from
the evaluation benchmark of Section~\ref{sec:dataset}, although both draw on the
same domain.

\paragraph{Vocabulary reuse.}
For a domain centered on people and organizations, established vocabularies such
as \texttt{foaf:Person} or \texttt{schema.org/Person} are natural reuse
candidates, and in an RDF setting we would align to them. We deliberately did not
import them here for two reasons. First, the deployment target processes sensitive
personnel records under data-sovereignty constraints and required Portuguese,
domain-specific labels (e.g., \texttt{patentePessoa} for military rank) with no
direct counterpart in general vocabularies. Second, the property-graph model we
adopted has no standardized OWL-style import mechanism, so reuse would be nominal
rather than semantic. We regard alignment to \texttt{foaf}/\texttt{schema.org}
(and publication of the ontology in OWL) as worthwhile future work for
interoperability, and we note it as a limitation in Section~\ref{sec:limitations}.

\paragraph{Why a property graph.}
We selected Neo4j and a labeled property graph---rather than RDF/SPARQL/OWL---for
pragmatic engineering reasons: native multi-hop traversal performance, mature
tooling and drivers for the local, on-premises deployment our use case demands,
and direct support from the \texttt{GraphCypherQAChain} translation tooling. The
cost of this choice, acknowledged here, is reduced access to formal reasoning and
to standardized validation (e.g., SHACL); Section~\ref{sec:consistency} shows that
a schema/ontology-aware validation step is precisely what the dominant error class
would benefit from, which makes the trade-off a substantive one rather than a
neutral implementation detail.

\paragraph{Resulting schema.}
The schema defines fourteen node labels: \texttt{Pessoa} (Person),
\texttt{Competencia} (Skill), \texttt{Projeto} (Project), \texttt{Conhecimento}
(Knowledge), \texttt{Tecnologia} (Technology), \texttt{Treinamento} (Training),
\texttt{Certificacao} (Certification), \texttt{Documentacao} (Documentation),
\texttt{Processo} (Process), \texttt{Pesquisa} (Research), \texttt{Publicacao}
(Publication), \texttt{Graduacao} (undergraduate degree), \texttt{PosGraduacao}
(graduate degree), and \texttt{Sumario} (abstract). Attributes such as a person's
rank and position are stored as \texttt{Pessoa} properties
(\texttt{patentePessoa}, \texttt{cargoPessoa}) rather than as separate nodes.
These labels are connected by eighteen typed, directed relationship types such as
\texttt{TEM\_COMPETENCIA} (has skill), \texttt{TRABALHA\_EM} (works in),
\texttt{UTILIZA\_TECNOLOGIA} (uses technology), \texttt{CONCLUIU\_TREINAMENTO}
(completed training), and \texttt{POSSUI\_CERTIFICACAO} (holds certification). The
schema is shown in Figure~\ref{fig:schema}. The graph was implemented in Neo4j~\cite{Neo4j2024} and
populated with synthetic but representative data using Python scripts
(Section~\ref{sec:dataset} discusses the synthetic data and its limitations). The
populated graph contains 145 nodes (70 \texttt{Pessoa}, the rest distributed over
the remaining labels) and 1{,}342 relationships across the eighteen types, the most
frequent being \texttt{INDICADO\_PARA} and \texttt{CONCLUIU\_TREINAMENTO} (215
each), \texttt{TRABALHA\_EM} (196), \texttt{TEM\_COMPETENCIA} (169), and
\texttt{POSSUI\_CONHECIMENTO} (158).

\begin{figure}[ht]
  \centering
  \includegraphics[width=0.80\linewidth]{Captura de tela 2026-04-29 203225.png}
  \caption{The Knowledge Graph schema used in this study. Nodes represent
           domain entity types; edges represent typed relationships.}
  \label{fig:schema}
\end{figure}

\subsection{GraphRAG Pipeline}
\label{sec:pipeline}

The pipeline has four stages, illustrated in Figure~\ref{fig:pipeline}:
(i)~a \textit{translation} module converts the natural language question into
Cypher using an LLM; (ii)~a \textit{retrieval} module executes that Cypher against
the KG and returns a precise result set; (iii)~an \textit{augmentation} module
serializes the result set as structured (JSON) context; and (iv)~a
\textit{generation} module produces a grounded natural language answer. Compared
to vector RAG, this eliminates retrieval approximation---if the Cypher is correct,
the facts are exact---but shifts the principal failure mode to query semantic
correctness. It was implemented in Python using Neo4j's official driver and
LangChain~\cite{LangChain2023} for prompt management.

\begin{figure}[ht]
  \centering
  \includegraphics[width=0.84\linewidth]{Captura de tela 2026-04-29 204216.png}
  \caption{The GraphRAG pipeline. The user submits a natural language question;
           an LLM translates it into Cypher; Neo4j executes the query and returns
           structured data; a second LLM step generates the final answer grounded
           in that data.}
  \label{fig:pipeline}
\end{figure}

\paragraph{Prompts and control flow.}
For reproducibility we make the prompt design explicit. The translation prompt
supplies the LLM with the full graph schema (node labels, relationship types, and
property keys) and the user question, and instructs it to return a single Cypher
statement. The generation prompt receives \emph{both the original natural-language
question and the JSON result set}, and is instructed to answer the question using
only the supplied data, without recourse to parametric knowledge. When the query
returns an empty result set, the generator is instructed to emit a fixed phrase
(``I did not find information about this in the database'') rather than
speculate.

\paragraph{Configurations.}
Three configurations were evaluated to disentangle the contributions of query
generation and answer generation:
\begin{itemize}
  \item \textbf{GPT-4 standalone.} GPT-4~\cite{openai2023gpt4} handles both the
        NL-to-Cypher translation and the final answer generation via the OpenAI
        API. This represents the cloud-based LLM-as-a-service paradigm and serves
        as the primary accuracy baseline.
  \item \textbf{HF\,27B standalone.} Both stages use the Hugging Face
        \texttt{text-to-cypher-Gemma-3-27B-Instruct} model, quantized to 4-bit
        precision (Q4) and served locally via Ollama. It was selected because it
        runs entirely on local hardware---a key constraint for organizations
        processing sensitive personnel data---and because it is specifically
        fine-tuned for Cypher generation, enabling a comparison between domain
        specialization and general capability.
  \item \textbf{Hybrid.} HF\,27B generates the Cypher query; GPT-4 summarizes the
        JSON result into natural language. Because the answer generator is held
        fixed (GPT-4) and only the upstream query source changes, the gap between
        the hybrid and GPT-4 standalone isolates the cost of the query-generation
        stage.
\end{itemize}
Local inference ran on a consumer workstation with an NVIDIA GeForce RTX\,5060\,Ti
(16\,GB GDDR6), AMD Ryzen~5 9600X, and 32\,GB DDR5 RAM.

\subsection{Evaluation Dataset and Metrics}
\label{sec:dataset}

Taking the critique of Zhang et al.~\cite{zhang2025diagnosingaddressingpitfallskgrag}
seriously, we applied a hybrid but primarily manual construction procedure. An
auxiliary LLM first generated a pool of candidate questions by systematically
varying entity types and relationship combinations present in the ontology. Every
candidate was then revised by the authors: questions were rewritten for clarity
and appropriate complexity; gold-standard Cypher queries were written and tested
directly in Neo4j rather than being LLM-generated; and reference answers were
composed manually to ensure factual correctness against the graph. The final
corpus contains 95 question--query--response triples.

\paragraph{Query-complexity breakdown.}
To characterize the benchmark---and to enable a per-class analysis
(Section~\ref{sec:consistency})---we classify each item by the structure of its
gold Cypher query into four classes: \emph{property lookups} (single node, no
traversal; 18 items), \emph{single-hop} relational queries (one relationship
pattern; 70 items), \emph{multi-pattern} conjunctive queries (two or more
patterns, e.g., ``uses both Python and PostgreSQL''; 6 items), and
\emph{aggregation} queries (1 item). This distribution is itself informative: the
benchmark is dominated by single-hop lookups, which we flag as a representativeness
limitation in Section~\ref{sec:limitations}. We also note that the questions,
while reviewed by the authors, were not validated by an external domain expert, and
that establishing their representativeness of real user queries is future work.

\paragraph{Metrics.}
Four metrics were computed for each configuration:
\begin{itemize}
  \item \textbf{Cypher Execution Rate.} The fraction of generated queries that
        Neo4j executes without error. This captures syntactic validity and
        \emph{schema compatibility at runtime}---i.e., that the query references
        only node labels, relationship types, and property keys that exist in the
        graph. It does \emph{not} certify that the query is semantically correct.
  \item \textbf{Factual Consistency (LLM-as-a-Judge).} Each generated response was
        rated \emph{Consistent}, \emph{Partially Consistent}, or \emph{Inconsistent}
        against the gold-standard answer, on a structured rubric, by an LLM judge.
        We use \textbf{Claude Opus~4 (Anthropic, model \texttt{claude-opus-4-8})}.
        Crucially, the judge is \emph{not} one of the three evaluated configurations,
        which removes the self-preference risk that arises when a system under test
        also judges. An answer is rated \emph{Consistent} when it conveys the gold
        answer's entities and facts without contradiction (extra grounded detail is
        allowed), \emph{Partially Consistent} when correct but incomplete, and
        \emph{Inconsistent} otherwise (wrong entities, fabrication, or a ``not
        found'' reply when data exist). This is our primary accuracy signal; a
        comparison against a GPT-4 judge is reported in Section~\ref{sec:consistency}.
  \item \textbf{BERTScore F1}~\cite{bert-score}. Token-level semantic similarity
        between the generated response and the gold standard, used as a secondary
        linguistic-quality indicator.
  \item \textbf{Cost and latency.} Wall-clock evaluation time and monetary cost per
        configuration (Section~\ref{sec:cost}).
\end{itemize}

\paragraph{Statistical reporting.}
With $n=95$ we report Consistent rates as rounded percentages with Wilson 95\%
confidence intervals~\cite{wilson1927}, avoiding two-decimal rates that would
overstate precision at this sample size. Paired configurations are compared on
categorical consistency with McNemar's exact test~\cite{mcnemar1947} and on
BERTScore with the Kruskal--Wallis and Wilcoxon signed-rank tests; agreement
between the two LLM judges is quantified with Cohen's $\kappa$~\cite{cohen1960}.
The script and labelled data that reproduce every reported figure are released
with the paper (Section~\ref{sec:limitations}).

%=======================================================================
\section{Results and Discussion}
%=======================================================================

\subsection{Infrastructure and Cost Trade-offs}
\label{sec:cost}

Before turning to accuracy, Table~\ref{tab:cost} records the practical costs of
each configuration, as these shape real deployment decisions. Evaluating the full
95-question dataset took approximately 25~minutes with GPT-4, at roughly
US\$\,7.50 in API fees. HF\,27B required around five hours on the available
hardware; its marginal monetary cost was negligible, as both inference and
consistency judging were performed offline. The hybrid configuration combines a local HF\,27B
translation pass (part of the same five-hour local budget) with a GPT-4
summarization pass of about 10~minutes and US\$\,4.70.

\begin{table}[ht]
  \caption{Cost and latency over the 95-question benchmark. Local time for the
           hybrid refers to the HF\,27B translation pass; cloud time refers to the
           GPT-4 summarization pass.}
  \label{tab:cost}
  \centering
  \begin{tabular}{lccc}
    \toprule
    \textbf{Configuration} & \textbf{Wall-clock time} & \textbf{Hardware} & \textbf{Monetary cost} \\
    \midrule
    GPT-4 standalone  & $\approx$25 min            & cloud API         & $\approx$US\$\,7.50 \\
    HF\,27B standalone & $\approx$5 h              & local GPU         & $\approx$US\$\,0\textsuperscript{*} \\
    Hybrid            & $\approx$5 h (local) + 10 min (cloud) & local GPU + cloud API & $\approx$US\$\,4.70 \\
    \bottomrule
  \end{tabular}
  \\[2pt]
  \footnotesize\textsuperscript{*}Local inference only; near-zero marginal monetary cost (consistency judging done offline).
\end{table}

These figures illustrate a trade-off that goes beyond accuracy: cloud-based APIs
offer low latency and high performance with no infrastructure investment, but
introduce recurring costs, vendor dependency, and potential data exposure; local
models invert this profile, offering data sovereignty and near-zero marginal cost
at the expense of hardware investment and significantly higher inference time. For
organizations handling sensitive personnel records, the local option may be
preferable even at a performance penalty.

\subsection{Cypher Query Generation}
\label{sec:cypher}

Table~\ref{tab:performance} summarizes the quantitative results. The HF\,27B
model, fine-tuned specifically for text-to-Cypher translation, achieved a perfect
100\% execution rate; GPT-4 reached 97.9\% (93/95). Manual inspection of GPT-4's
two failures found that both originated from surface-level string formatting
issues---mismatched quote characters and missing label backticks---reported
uniformly by Neo4j as \texttt{Neo.ClientError.Statement.SyntaxError}, rather than
from any misunderstanding of the question. GPT-4 understood \emph{what} to retrieve
but occasionally stumbled on exact syntax, whereas HF\,27B's fine-tuning made it
essentially immune to these surface errors.

We note a low-cost remedy that would likely close this small gap entirely: feeding
the Neo4j syntax error back to GPT-4 with a request to repair the query. Because
this would be triggered in only 2.1\% of cases, the overhead is negligible; we did
not apply it here in order to keep the three configurations directly comparable,
but it is the obvious first step for a production deployment and would reduce the
text-to-Cypher choice to ``cloud vs.\ fully local'' rather than requiring a
specialized model for syntactic robustness alone.

\begin{table*}[ht]
  \caption{Comparative performance of the three architectures on the 95-question
           gold-standard benchmark. The consistency categories and the
           error category together partition all 95 items in each column (the
           underlying counts sum to 95; rounded percentages may differ from 100 by
           up to one point). Factual consistency is judged by Claude Opus~4
           (\texttt{claude-opus-4-8}); bracketed values for the \emph{Consistent}
           row are Wilson 95\% confidence intervals.}
  \label{tab:performance}
  \centering
  \begin{tabular}{lccc}
    \toprule
    \textbf{Metric}
      & \textbf{GPT-4}
      & \textbf{HF\,27B}
      & \textbf{Hybrid (HF\,27B + GPT-4)} \\
    \midrule
    Cypher Execution Rate              & 97.9\% & 100\% & 100\% \\
    Query-level entity recall (vs.\ gold) & 63\% & 68\% & 68\% \\
    BERTScore F1 (mean)                & 0.793  & 0.760 & 0.742 \\
    \midrule
    Factual Consistency: Consistent    & 70\% [61--79] & 49\% [40--59] & 59\% [49--68] \\
    Factual Consistency: Partial       &  6\% & 6\%  &  9\% \\
    Factual Consistency: Inconsistent  & 21\% & 44\% & 32\% \\
    No answer (API/exec.\ error)\textsuperscript{\dag} &  2\% & 0\%  &  0\% \\
    \bottomrule
  \end{tabular}
  \\[2pt]
  \footnotesize\textsuperscript{\dag}Two GPT-4 items produced no answer because the
  generation call hit an API rate limit; with no text to assess they are counted as
  not consistent and listed here for transparency. With Claude as judge every item
  that has an answer was scored, so there are no judge-scoring errors (the earlier
  GPT-4 judge had left 6 HF\,27B and 2 hybrid items unscored).
\end{table*}

\subsection{Factual Consistency of Generated Answers}
\label{sec:consistency}

The accuracy picture reverses when we look at response quality. Judged by Claude
(Section~\ref{sec:dataset}), GPT-4 standalone produced consistent answers 70\% of
the time (95\% CI 61--79), the highest of the three; HF\,27B reached 49\%
(95\% CI 40--59) and the hybrid 59\% (95\% CI 49--68).

\paragraph{Both pipeline stages contribute to the gap.}
Two paired comparisons localize where quality is lost, both scored by the same
judge. Holding the \emph{query} fixed and changing only the generator (HF\,27B
standalone $\to$ hybrid) raised consistency from 49\% to 59\%, isolating the value
of a stronger summarizer (McNemar exact test, $p=0.06$). Holding the \emph{generator}
fixed at GPT-4 and changing only the query source (hybrid $\to$ GPT-4 standalone)
raised consistency from 59\% to 70\%, isolating the cost of HF\,27B's retrieval
(McNemar, $p=0.06$). The two stages thus contribute comparably, and at $n=95$ each
individual gap is only marginally significant; what is unambiguous is the
end-to-end difference between the fully local and the fully cloud pipeline
(HF\,27B 49\% vs.\ GPT-4 70\%, $p<0.001$). Neither stage can be skipped: a
semantically wrong query cannot be repaired by the generator, and verbose
whole-node retrieval defeats a weak summarizer. Because the judge (Claude) is not
one of the evaluated systems, none of these comparisons is affected by judge
self-preference. We examine the two failure routes next.

\paragraph{A direct query-level check.}
Because gold Cypher exists for all 95 items, we executed every generated query
against the populated graph and compared its result with the gold query's. Exact
result-set identity is low and---misleadingly---similar across configurations
(GPT-4 14\%, HF\,27B/hybrid 18\%): it is dominated by \emph{projection}
differences (selecting different properties, or returning a whole node where the
gold projects two fields), not by retrieving the wrong entities. A
projection-robust comparison that matches on the set of returned entity
\emph{names} gives a fairer picture: the generated query recovers all gold entities
in 63\% of items for GPT-4 and 68\% for HF\,27B (the hybrid inherits HF\,27B's
queries, hence 68\%), with exact name-set agreement of 61\% and 58\%,
respectively. The key observation is that HF\,27B's queries are \emph{at least as
accurate as GPT-4's at the entity level}---consistent with only 7/95 targeting the
wrong labels---so HF\,27B's far lower end-to-end consistency (49\%) is not caused
by worse retrieval but by the answer-generation step. This is the clearest evidence
that the local model's binding weakness is the summarization of its verbose,
whole-node results rather than query semantics, and it is why our primary accuracy
signal is the Claude-judged answer consistency rather than a raw result-set match.

\paragraph{What actually goes wrong upstream.}
A direct inspection of HF\,27B's 95 generated queries refines the usual ``the model
writes semantically wrong Cypher'' narrative. Genuine schema-element errors---where
the query targets a different node label than the gold query---occur in only
\textbf{7 of 95} cases. A representative example is question Q38, ``which people
are data scientists and master Neo4j?'': the gold query treats ``data scientist''
as a \texttt{Cargo} (position) attribute, whereas HF\,27B modelled it as a
\texttt{Competencia} (skill) and returned nothing
(Listing~\ref{lst:semantic}). This is the failure mode that matters
most---it executes cleanly and is invisible to the execution-rate
metric---but it is not the most frequent one.

\begin{lstlisting}[style=cypher,caption={A genuine semantic query error (Q38). The
gold query filters on the \texttt{Cargo} ``Cientista de Dados''; HF\,27B instead
looks for a \texttt{Competencia} containing ``ciência de dados'', which does not
exist, so the executable query returns an empty set.},label={lst:semantic}]
-- Gold (correct schema element: Cargo)
MATCH (p:Pessoa {cargoPessoa:'Cientista de Dados'})
      -[:DOMINA_TECNOLOGIA]->(t:Tecnologia {nomeTecnologia:'Neo4j'})
RETURN p.idPessoa, p.nomePessoa

-- HF 27B (wrong schema element: Competencia)
MATCH (p:Pessoa)-[:TEM_COMPETENCIA]->(c:Competencia)
WHERE toLower(c.nomeCompetencia) CONTAINS 'ciência de dados'
  AND (p)-[:DOMINA_TECNOLOGIA]->(:Tecnologia {nomeTecnologia:'Neo4j'})
RETURN p
\end{lstlisting}

More frequently, the query is schema-correct but the \emph{answer generation}
fails. HF\,27B returned the whole matched node (\texttt{RETURN p}) in 47 of 95
queries---versus 10 of 95 for GPT-4, which tends to project specific properties
(\texttt{RETURN p.nomePessoa}). Returning whole nodes produces a verbose, deeply
nested JSON payload that the summarization step handles poorly. Of HF\,27B's 51
``I did not find information'' answers, \textbf{46 were produced over a
schema-correct query}; the empties are therefore predominantly summarization
failures, not retrieval failures. The decisive evidence is what happens when the
\emph{same} retrieval is summarized by GPT-4 instead: in 11 of those 46 cases the
hybrid recovered a correct, non-empty answer. Question Q1 is the canonical example
(Listing~\ref{lst:summ}): given identical Cypher and identical results, HF\,27B
reported ``not found'' while GPT-4 correctly named the two people.

\begin{lstlisting}[style=cypher,caption={A summarization failure (Q1). HF\,27B and
the hybrid execute the \emph{same} query; HF\,27B's own generator reports an empty
answer, whereas GPT-4 summarizing the identical result set returns the correct
people. This isolates answer generation, not retrieval, as the fault.},label={lst:summ}]
-- Query generated by HF 27B (used by both HF 27B standalone and Hybrid)
MATCH (p:Pessoa)-[:TEM_COMPETENCIA]->(c:Competencia)
WHERE toLower(p.patentePessoa) CONTAINS 'major'
  AND toLower(c.nomeCompetencia) CONTAINS 'gestão de projetos'
RETURN p

-- HF 27B answer:  "I did not find information about this in the database."
-- Hybrid answer (GPT-4 over the same results): "Militar 62 and Militar 65 ..."
\end{lstlisting}

These three patterns---genuine schema-element errors (rare but consequential),
whole-node retrieval that overwhelms summarization, and prompt-instruction
violations (e.g., speculating instead of emitting the standardized empty-result
phrase)---are consistent with HF\,27B being tuned narrowly for text-to-Cypher
\emph{translation} rather than for the distinct task of summarizing structured
results. A different locally-run model dedicated to the generation step, or a
generation-aware fine-tune, is a natural remedy.

\paragraph{An ontology-aware validation step would help.}
The most consequential error---an executable query that retrieves the wrong
data---is exactly what schema- or ontology-aware validation could catch. Rather
than checking only that the query executes, the system could verify that it uses
the intended classes and relationships and that its result shape matches the
question's expected interpretation, via constrained/schema-guided generation, query
templates, or SHACL-style checks on the result set before answer generation. This
is the main semantic lever left unexploited by the current property-graph
implementation and a primary direction for future work.

\paragraph{Per-class breakdown.}
Table~\ref{tab:byclass} reports consistency by query class, scored by the same
Claude judge used throughout, so it is directly comparable with
Table~\ref{tab:performance}. The GPT-4 advantage is
concentrated in single-hop queries (73\% vs.\ 53\% for the hybrid and 44\% for
HF\,27B), where HF\,27B's whole-node retrieval most often derails summarization;
property lookups, by contrast, are handled comparably by all three, and the small
multi-pattern set is essentially even.

\begin{table}[ht]
  \caption{Consistent-answer rate by query class, scored by the same Claude judge
           used throughout, directly comparable with Table~\ref{tab:performance}. The
           benchmark is dominated by single-hop queries, where the upstream-query
           effect is largest.}
  \label{tab:byclass}
  \centering
  \begin{tabular}{lcccc}
    \toprule
    \textbf{Query class} & \textbf{$n$} & \textbf{GPT-4} & \textbf{HF\,27B} & \textbf{Hybrid} \\
    \midrule
    Property lookup  & 18 & 56\%  & 61\% & 67\%  \\
    Single-hop       & 70 & 73\%  & 44\% & 53\%  \\
    Multi-pattern    &  6 & 83\%  & 83\% & 100\% \\
    Aggregation      &  1 & 100\% & 0\%  & 100\% \\
    \bottomrule
  \end{tabular}
\end{table}

\paragraph{Judge reliability.}
Using an LLM as judge is defensible here for two reasons. First, our judge (Claude)
is not among the three evaluated configurations, so the self-preference bias that
arises when a system under test also judges is structurally absent. Second, to check
that the judgments are not idiosyncratic, we re-scored the items with a second,
independent LLM judge (GPT-4) and found high agreement with Claude---98\% exact
agreement (Cohen's $\kappa=0.95$) on the GPT-4 configuration and 94\% ($\kappa=0.89$)
on the hybrid, the two configurations for which per-item GPT-4-judge labels are
available. The two judges agree on the ranking and on the wide GPT-4-vs-HF\,27B gap
($p<0.001$ under both); they differ mainly in leniency toward verbose-but-grounded
answers, which slightly shifts absolute rates (the GPT-4 judge scored the hybrid at
53\% vs.\ Claude's 59\%). We adopt Claude as the primary judge because it is
bias-free by construction and because the GPT-4-judge run left eight items unscored
due to API rate-limiting. A human-annotated subsample remains the most desirable
next step.

\subsection{BERTScore Analysis}

BERTScore~\cite{bert-score} compares the contextual token embeddings of the
generated and reference texts, rewarding semantic equivalence even when surface
words differ (e.g., ``senior engineer'' vs.\ ``level-3 developer''). It is well
suited to crediting valid paraphrases of a fact, but---because it operates on
token-embedding similarity---it does \emph{not} verify entity identity or numeric
precision: an answer naming the wrong person in the right vocabulary still scores
fairly high. We therefore treat factual consistency as primary and BERTScore as a
secondary, complementary signal.

Figure~\ref{fig:bert_box} shows the BERTScore F1 distribution for the three
configurations as box plots, which summarize the median, quartiles, and spread far
more legibly than the per-question scatter plots of the previous version. GPT-4 has
the highest median (0.82) and the two near-zero outliers correspond exactly to its
two execution failures. HF\,27B (mean 0.760) is more dispersed with a longer lower
tail, reflecting summarization over malformed or empty context; the hybrid (mean
0.742) has the lowest median. A Kruskal--Wallis test confirms the distributions
differ ($H=10.8$, $p=0.004$); pairwise Wilcoxon signed-rank tests show GPT-4
significantly above both HF\,27B ($p=0.002$) and the hybrid ($p<0.001$), while
HF\,27B vs.\ hybrid is \emph{not} significant ($p=0.08$).

\begin{figure}[ht]
  \centering
  \includegraphics[width=0.85\linewidth]{bertscore_boxplot.png}
  \caption{BERTScore F1 distribution per configuration ($n=95$). Boxes show
           the interquartile range and median; diamonds mark the mean; points are
           individual questions (jittered). GPT-4's two near-zero points are its
           execution failures.}
  \label{fig:bert_box}
\end{figure}

The counter-intuitive result---that the hybrid scores no higher than HF\,27B on
BERTScore despite using the stronger GPT-4 generator---has a concrete explanation
that the non-significant difference already tempers. Because HF\,27B returns whole
nodes, GPT-4 in the hybrid is handed verbose JSON and tends to produce longer,
detail-padded answers (e.g., in Q1 the hybrid appended each person's position,
admission date, and e-mail, earning a \emph{Partially Consistent} rating). These
additions are factually grounded but diverge in surface form from the concise,
manually written gold answers, which depresses token-level similarity. This is an
interaction between retrieval shape and generation, not evidence that GPT-4 is a
weaker summarizer; it also illustrates why we do not lean on BERTScore as a primary
metric for a factual-traceability setting.

\subsection{Application Scenarios}

The benchmark queries reflect genuine organizational information needs. We
highlight four areas where the architecture is immediately actionable.
\textbf{Strategic workforce planning}: queries such as \textit{``list all
technologies required by Project~X for which we have fewer than three practitioners
at senior level''} expose skill gaps for proactive hiring and training.
\textbf{Knowledge succession}: \textit{``which knowledge domains critical to
Project~Y are held by fewer than two people?''} surfaces institutional fragility
before it becomes a crisis---directly relevant in organizations with high turnover
or mandatory rotation, such as the military context that motivated this study.
\textbf{Dynamic personnel allocation}: \textit{``find all active specialists
holding certification~Z who completed training~W in the last 24 months and are not
on a critical project''} enables rapid, auditable staffing decisions---the answer
is a traversal result, not an LLM opinion. \textbf{Legal and regulatory research}:
a KG of legislation, decisions, and citations would let compliance officers ask
provenance questions whose answers are grounded in graph structure, where
hallucination carries direct legal consequence.

%=======================================================================
\section{Limitations}
\label{sec:limitations}
%=======================================================================

Several limitations bound our conclusions. \textbf{Synthetic data and benchmark
scope}: the KG was populated with synthetic data and the 95-item benchmark, while
manually validated, was seeded by an LLM and not reviewed by an external domain
expert; it is also skewed toward single-hop queries (Section~\ref{sec:dataset}), so
the results may not transfer to richer multi-hop or aggregation workloads. A real
organizational case study is needed. \textbf{Statistical power}: with $n=95$,
confidence intervals are wide and some differences (e.g., HF\,27B vs.\ hybrid on
BERTScore) are not significant; absolute rates should be read with their intervals.
\textbf{Judge}: factual consistency is assessed by an LLM judge (Claude,
\texttt{claude-opus-4-8}) deliberately not among the evaluated systems; a second LLM
judge (GPT-4) agrees closely (98\%/94\%, $\kappa\geq0.89$), but no human annotation
was performed, which remains desirable. \textbf{Single-run estimates}: each configuration was evaluated
once, without multiple seeds or temperature settings, so run-to-run variance is not
quantified. \textbf{Unexploited semantics}: the KG is used as a retrieval store
without reasoning, constraints, or validation, and the ontology is not yet
published in OWL or aligned to standard vocabularies.

\paragraph{Reproducibility.}
To address reproducibility, the domain ontology, the synthetic-data generation
scripts, the Neo4j dump, the 95-item benchmark, the pipeline prompts, the labelled
per-item evaluation data, and the analysis script (\texttt{analysis.py}) that
recomputes every reported statistic are made available at
\url{https://github.com/antgustavo/Management-Knowledge}.\footnote{Repository URL to be
inserted for the camera-ready version.}

%=======================================================================
\section{Conclusion}
%=======================================================================

This paper evaluated a GraphRAG architecture for strategic personnel management
across three LLM configurations and a 95-question gold-standard benchmark. The
central finding is a clear division of labor: a domain-specialized model excels at
syntactic Cypher generation (100\% execution for HF\,27B), while a general-purpose
proprietary model leads on factual consistency of the generated answer (70\% for
GPT-4, judged by Claude). Decomposing the local-vs-cloud gap with a single fixed
judge shows that the query-generation and answer-generation stages contribute
comparably (each isolated step raises consistency by about ten points, each only
marginal at $n=95$), while the end-to-end gap is large and robust (49\% vs.\ 70\%,
$p<0.001$). A finer-grained analysis shows that HF\,27B's deficit is driven less by
frankly wrong queries (7/95) than by schema-correct queries whose whole-node result
shape defeats summarization---a distinction with direct design implications.

Beyond accuracy, the comparison surfaces a cloud-versus-local trade-off relevant to
real organizations: GPT-4 offers the highest accuracy with minimal infrastructure
but recurring cost and external data exposure; HF\,27B offers data sovereignty and
near-zero marginal cost but currently lags on end-to-end answer quality. The path
forward is concrete: (i)~add an ontology-aware query-validation step
(constrained generation, templates, or SHACL-style result checks) to catch the
executable-but-wrong queries; (ii)~fine-tune an open model end-to-end on full
pipeline trajectories (question~$\rightarrow$~Cypher~$\rightarrow$~JSON~$\rightarrow$~answer)
and project specific properties rather than whole nodes; (iii)~evaluate larger,
unquantized models on server-class hardware to separate quantization from capacity;
and (iv)~build larger, expert-validated KG-RAG benchmarks following the criteria of
Zhang et al.~\cite{zhang2025diagnosingaddressingpitfallskgrag} to enable
statistically robust comparison.

\begin{acknowledgments}
We thank the Instituto Tecnológico de Aeronáutica (ITA) for institutional
support. We also thank Major Aviador Douglas Silva Teixeira (CCA-SJ) for
co-advising the undergraduate thesis from which this paper originates, and the
open-source communities behind Neo4j, LangChain, and the Hugging Face ecosystem
for making this research reproducible.
\end{acknowledgments}

\section*{Declaration on Generative AI}
During the preparation of this work the authors used Google Gemini to assist
with English translation and \LaTeX{} formatting. All content was subsequently
reviewed and edited by the authors, who take full responsibility for the final
publication.

\bibliography{sample-ceur}

\end{document}
